% ==========================================
% Step 4/4 - Final refinement (refinement step 3)
% Output: step4-2026-10-15-wednesday-week5-algebra-1-offset-final.tex
% Input: step3-2026-10-15-wednesday-week5-algebra-1-offset-speech_refined.tex
% Model: gemini-3.8-flash
% Temperature: 0.35
% TopP: 0.9
% TopK: 10
% MaxOutputTokens: 65535
% ThinkingBudget: 4096
% ThinkingLevel: MEDIUM
% Processed on: 2026-10-17 00:18:10
% Thought summaries: step4-2026-10-15-wednesday-week5-algebra-1-offset-final-thoughts.md
% ------------------------------------------
% Token Usage (1 Request(s)):
%   - Total Prompt Tokens : 57’121 (gesamter Prompt, inkl. Cache)
%   - Cached Context      : 0 (aus Google Context-Cache, rabattiert)
%   - Fresh Input Tokens  : 57’121 (neu verrechneter Input)
%   - Output Tokens       : 25’061 (generiertes LaTeX, ohne Thinking)
%   - Thinking Tokens     : 10’873 (Denk-Tokens, als Output verrechnet)
% ==========================================
% ==========================================
% Step 3/4 - Speech refinement (refinement step 2)
% Output: step3-2026-10-15-wednesday-week5-algebra-1-offset-speech_refined.tex
% Input: step2-2026-10-15-wednesday-week5-algebra-1-offset-merged.tex
% Model: gemini-3.8-flash
% Temperature: 0.35
% TopP: 0.9
% TopK: 10
% MaxOutputTokens: 65535
% ThinkingBudget: 4096
% ThinkingLevel: MEDIUM
% Processed on: 2026-10-17 00:11:53
% Thought summaries: step3-2026-10-15-wednesday-week5-algebra-1-offset-speech_refined-thoughts.md
% ------------------------------------------
% Token Usage (1 Request(s)):
%   - Total Prompt Tokens : 149’973 (gesamter Prompt, inkl. Cache)
%   - Cached Context      : 0 (aus Google Context-Cache, rabattiert)
%   - Fresh Input Tokens  : 149’973 (neu verrechneter Input)
%   - Output Tokens       : 24’568 (generiertes LaTeX, ohne Thinking)
%   - Thinking Tokens     : 12’763 (Denk-Tokens, als Output verrechnet)
% ==========================================
% ==========================================
% Step 2/4 - Merge & timestamp alignment (refinement step 1)
% Output: step2-2026-10-15-wednesday-week5-algebra-1-offset-merged.tex
% Input: step1-2026-10-15-wednesday-week5-algebra-1-all-offset.tex
% Model: gemini-3.8-flash
% Temperature: 0.4
% TopP: 0.95
% TopK: 40
% MaxOutputTokens: 65535
% ThinkingBudget: 32765
% ThinkingLevel: MEDIUM
% Processed on: 2026-10-17 00:06:24
% Thought summaries: step2-2026-10-15-wednesday-week5-algebra-1-offset-merged-thoughts.md
% ------------------------------------------
% Token Usage (1 Request(s)):
%   - Total Prompt Tokens : 111’253 (gesamter Prompt, inkl. Cache)
%   - Cached Context      : 0 (aus Google Context-Cache, rabattiert)
%   - Fresh Input Tokens  : 111’253 (neu verrechneter Input)
%   - Output Tokens       : 24’858 (generiertes LaTeX, ohne Thinking)
%   - Thinking Tokens     : 10’505 (Denk-Tokens, als Output verrechnet)
% ==========================================
% ==========================================
% Step 1/4 - Transcription, 2 part(s) combined
% Output: step1-2026-10-15-wednesday-week5-algebra-1-all-offset.tex
% Source video: 2026-10-15-wednesday-week5-algebra-1.mp4
% Model: gemini-3.8-flash
% Temperature: 0.36
% TopP: 0.95
% TopK: 40
% MaxOutputTokens: 65535
% ThinkingBudget: 32765
% ThinkingLevel: HIGH
% Processed on: 2026-10-16 23:39:15
% ------------------------------------------
% Parts:
%   - Part 1: 2026-10-15-wednesday-week5-algebra-1-part1.mp4 | start 00:00:00 | model gemini-3.8-flash | processed 2026-10-16 23:30:40 | prompt 342’052 | output 13’964 | thinking 31’270
%   - Part 2: 2026-10-15-wednesday-week5-algebra-1-part2.mp4 | start 00:19:01 | model gemini-3.8-flash | processed 2026-10-16 23:39:15 | prompt 355’684 | output 14’088 | thinking 25’646
% ------------------------------------------
% Token Usage of this run (2 Request(s)):
%   - Total Prompt Tokens : 697’736 (gesamter Prompt, inkl. Cache)
%   - Cached Context      : 0 (aus Google Context-Cache, rabattiert)
%   - Fresh Input Tokens  : 697’736 (neu verrechneter Input)
%   - Output Tokens       : 28’052 (generiertes LaTeX, ohne Thinking)
%   - Thinking Tokens     : 56’916 (Denk-Tokens, als Output verrechnet)
% ==========================================

% Primary Language: English
% End of the video: 00:42:03

\lecturechapter[5]{Wednesday}{15 Oct}{15 October 2026}{The Chinese Remainder Theorem and Polynomial Rings}
% Old chapter title: The Chinese Remainder Theorem and Polynomial Rings

\begin{overview}[Course Content Overview]
In this lecture, we complete the general ideal-theoretic formulation of the Chinese Remainder Theorem for commutative rings and begin the theory of polynomial rings over fields:
\begin{itemize}
    \item \textbf{The Canonical Product Homomorphism:} Definition of the canonical map $\phi \colon \faktor{R}{IJ} \to \faktor{R}{I} \times \faktor{R}{J}$ and verification of well-definedness via the ideal inclusions $IJ \subseteq I, J$.
    \item \textbf{Failure of Isomorphism for Non-Coprime Ideals:} A concrete counterexample in $\mathbb{Z}$ with $I = (4)$ and $J = (6)$, illustrating that $\phi$ fails to be either injective or surjective when ideals are not coprime.
    \item \textbf{The Chinese Remainder Theorem:} Proof that comaximality ($I + J = (1)$) implies that $\phi$ is an isomorphism, utilizing the partition of unity identity $i + j = 1$.
    \item \textbf{Applications to $\mathbb{Z}$ and Euler's Totient Function:} Iterated factorization into prime powers, the multiplicative structure of unit groups, and the derivation of Euler's product formula for $\varphi(n)$ alongside the irregular prime $691$ and Bernoulli numbers.
    \item \textbf{Polynomial Rings over Fields:} Rigorous definition of the degree function on $K[x]$, the convention of leaving $\deg(0)$ undefined, and polynomial long division with remainder.
    \item \textbf{Principal Ideal Domains (PIDs):} Proof via well-ordering of degrees that every ideal in $K[x]$ is principal, the formal definition of a PID, and non-examples such as $\mathbb{Z}[x]$ and $\mathbb{C}[x, y]$.
\end{itemize}
\end{overview}

\section{The Chinese Remainder Theorem for Commutative Rings}

\subsection{Ideals and the Canonical Homomorphism}

\begin{speech}[00:00:00 - 00:02:30]
Okay, so, we start. And yesterday I was in the middle of this discussion of the Chinese Remainder Theorem. \inlinemetanote{writes: Chinese RT:}

And I will complete it now, but we start from a slightly different perspective. Let's just start with $I$ and $J$ \inlinemetanote{writes: $I, J$}... Isn't that loud? Doesn't it seem loud? No? I think that's somehow louder than usual.

Okay, so, we start with $I$ and $J$ inside a ring $R$ \inlinemetanote{writes: $\subset R$}, so, these are ideals \inlinemetanote{writes: ideals}. And I had discussed... ah, thanks. \inlinemetanote{student brings chalk} Thanks.

I had discussed last time that there's a notion of adding ideals \inlinemetanote{writes: $I+J$} and multiplying ideals \inlinemetanote{writes: $IJ$}. So, these are both ideals of $R$ \inlinemetanote{writes: ideals of $R$}.

And without... just assuming just this thing: we take $I$ and $J$ are ideals of $R$, we can construct this quotient $R \bmod IJ$ \inlinemetanote{writes: $\faktor{R}{IJ}$}, and we could construct $R \bmod I \times R \bmod J$ \inlinemetanote{writes: $\faktor{R}{I} \times \faktor{R}{J}$}. And moreover, we can construct this homomorphism $\phi$ \inlinemetanote{writes: $\phi \colon \faktor{R}{IJ} \to \faktor{R}{I} \times \faktor{R}{J}$}, and the rule for $\phi$ is $\phi(\bar{z}) = (\bar{z}, \bar{z})$ \inlinemetanote{writes: $\phi(\bar{z}) = (\bar{z}, \bar{z})$}. So, it's kind of almost a silly, tautological formula.

And the interesting thing is you have to remember that this $\bar{z}$ is... the bar here is $\bmod IJ$ \inlinemetanote{writes \qt{$\bmod IJ$} under the first $\bar{z}$}, and the bar here is $\bmod I$ \inlinemetanote{writes \qt{$\bmod I$} under the second $\bar{z}$}, and the bar here means $\bmod J$ \inlinemetanote{writes \qt{$\bmod J$} under the third $\bar{z}$}. So, the three bars have three different meanings, which you have to be able to read from the place they appear in this formula.

And this is always a homomorphism: $\phi$ is a homomorphism in general \inlinemetanote{writes: $\phi$ is a hom.}. A ring homomorphism in general. This part of the theory just holds in arbitrary generality. Just because the only thing we had to check here was that this definition works, and it works because this is well-defined \inlinemetanote{writes: Well defined} --- I mentioned this, but I'll say it explicitly ---, well-defined because $IJ$ is inside $I$ \inlinemetanote{writes: $IJ \subset I$} and $IJ$ is inside $J$ \inlinemetanote{writes: $IJ \subset J$}. That makes this association well-defined. All right?
\end{speech}

\recap{Sums and products of ideals (\cref{def:sum-of-ideals}, \cref{def:product-of-ideals}), the canonical map $\phi$ (\cref{def:crt-canonical-homomorphism}) with its well-definedness (\cref{lem:phi-well-defined}), and the statement of the Chinese Remainder Theorem (\cref{thm:crt-general-rings}) are from the previous lecture. The lecturer recalls them and proves the theorem today.}

\begin{content}[The Canonical Product Homomorphism]
Let $R$ be a commutative ring with identity, and let $I, J \subseteq R$ be ideals.
\begin{itemize}
    \item \textbf{Sum of Ideals:} $I + J := \{ i + j \mid i \in I, \, j \in J \}$ is the smallest ideal of $R$ containing both $I$ and $J$.
    \item \textbf{Product of Ideals:} $IJ := \left\{ \sum_{k=1}^m i_k j_k \;\middle|\; m \in \mathbb{N}, \, i_k \in I, \, j_k \in J \right\}$ is an ideal of $R$ satisfying $IJ \subseteq I \cap J \subseteq I, J$.
\end{itemize}

\begin{definition}[Canonical Product Homomorphism]\label[definition]{def:canonical-crt-map}
The canonical map between quotient rings is defined by:
\begin{equation}\label{eq:canonical-crt-map}
\begin{aligned}
\phi \colon \faktor{R}{IJ} &\longrightarrow \faktor{R}{I} \times \faktor{R}{J}, \\
z + IJ &\longmapsto (z + I, \, z + J),
\end{aligned}
\end{equation}
which on the blackboard is written compactly with class bars as:
\[
\phi(\underbrace{\bar{z}}_{\bmod IJ}) = (\underbrace{\bar{z}}_{\bmod I}, \, \underbrace{\bar{z}}_{\bmod J}).
\]
\end{definition}

\begin{proposition}[Well-Definedness and Homomorphism Property]\label[proposition]{prop:phi-well-defined}
The mapping $\phi$ is a well-defined homomorphism of rings with identity.
\end{proposition}

\begin{proof}
Let $z, z' \in R$ with $\bar{z} = \bar{z}'$ in $\faktor{R}{IJ}$, meaning $z - z' \in IJ$. Since $IJ \subseteq I$ and $IJ \subseteq J$, we immediately have:
\[
z - z' \in I \implies z + I = z' + I \quad \text{and} \quad z - z' \in J \implies z + J = z' + J.
\]
Hence $\phi(z + IJ) = \phi(z' + IJ)$, proving well-definedness. That $\phi$ preserves addition, multiplication, and the identity $(1 + I, 1 + J)$ follows directly from the componentwise definition of ring operations in the direct product.
\end{proof}

\begin{steps}[Semantic Meaning of the Class Bars]
Although the symbol $\bar{z}$ appears on both sides of $\phi(\bar{z}) = (\bar{z}, \bar{z})$, the three occurrences denote equivalence classes in three distinct quotient rings:
\begin{enumerate}[label=\textbf{(\arabic*)}]
    \setcounter{enumi}{0} \item \textbf{Domain Class:} The argument $\bar{z}$ represents the coset $z + IJ \in \faktor{R}{IJ}$.
    \setcounter{enumi}{1} \item \textbf{First Component:} The first component $\bar{z}$ represents the coset $z + I \in \faktor{R}{I}$.
    \setcounter{enumi}{2} \item \textbf{Second Component:} The second component $\bar{z}$ represents the coset $z + J \in \faktor{R}{J}$.
\end{enumerate}
\end{steps}
\end{content}

\subsection{A Counterexample to Isomorphism in \texorpdfstring{$\mathbb{Z}$}{Z}}

\begin{speech}[00:02:30 - 00:04:48]
And so, let's do an example before we go to the Chinese Remainder Theorem. Example is, let's say, we have the ideal $(4)$ \inlinemetanote{writes: $(4) = I$}, that's equal to $I$, and the ideal $(6)$ that's equal to $J$ \inlinemetanote{writes: $(6) = J$}. And these are ideals inside, of course, our favorite ring, the integers \inlinemetanote{writes: $\subset \mathbb{Z}$}. Well, maybe it's your favorite ring, I don't know. It's not my favorite ring, but... okay.

So, then what we claim here is that we have $\mathbb{Z}$... and now this is the product ideal, and we discussed last time that the product ideal here is just $(24)$; it's the product \inlinemetanote{writes: $\faktor{\mathbb{Z}}{(24)} = \faktor{\mathbb{Z}}{(4)(6)}$}. That's equal to the product ideal $(4)(6)$ \inlinemetanote{writes: $(4)(6)$}, okay? Do you agree?

And this thing maps, by that map on the board, to $\mathbb{Z} \bmod 4$ --- so what's this for? --- and $\mathbb{Z} \bmod 6$ \inlinemetanote{writes: $\xrightarrow{\phi} \faktor{\mathbb{Z}}{(4)} \times \faktor{\mathbb{Z}}{(6)}$}. Sorry. So, sometimes we write this $(4)$ as $4\mathbb{Z}$, it just means the same thing; it depends on your mood. Here I write it in this... these are different ways of writing an ideal, right? Today we're writing it like that.

Okay, so, this is the map. So, first of all, is $\phi$ surjective? \inlinemetanote{writes: Is $\phi$ surjective?} And I'm just going to answer this for you: the answer is going to be no \inlinemetanote{writes: No.}.

And why? For example, I could consider the point here that's given, so, consider $(0, 1)$ \inlinemetanote{writes: consider $(0, 1)$}. All right, that's a point, $0$ and $1$. Could that come from anything? Well, it would have to come from something divisible by $4$, so, it would have to be even. But then it's $1 \bmod 6$, so, it would have to be odd. It's very hard for it to be both even and odd; it's not possible, in fact. So, this is not in the image \inlinemetanote{writes: $(0, 1) \notin \operatorname{Im}(\phi)$}. Are you happy with that argument?

So, it's not surjective. Is it injective? \inlinemetanote{writes: Is $\phi$ inj.?} And the answer is going to be also no. But can you tell me an element here that goes to zero? An element that's not zero that goes to zero? Yeah?
\end{speech}

\begin{student}[Student Answer]
12.
\end{student}

\begin{speech}[00:04:51 - 00:05:25]
Well, zero is zero, so $12$ is... So, if you take here $\phi(\bar{12})$ \inlinemetanote{writes: $\phi(\bar{12})$}, this is equal to $(\bar{0}, \bar{0})$ \inlinemetanote{writes: $= (\bar{0}, \bar{0})$}. So, that shows that $\bar{12}$ goes to zero. And it's not zero because $24 \nmid 12$.

So, it's neither injective nor surjective. That's the typical thing. If you pick some ring with these ideals, you can construct this map, it looks pretty good, it's a homomorphism, but in general it's neither injective, nor is it surjective. And so, what the Chinese Remainder Theorem says...
\end{speech}

\begin{content}[Failure of Isomorphism for Non-Coprime Ideals]
\begin{example}[The Homomorphism for $I = (4)$ and $J = (6)$ in $\mathbb{Z}$]\label[example]{ex:non-coprime-crt}
In the ring of integers $R = \mathbb{Z}$, consider the principal ideals $I = (4) = 4\mathbb{Z}$ and $J = (6) = 6\mathbb{Z}$.
Their product ideal is:
\[
IJ = (4)(6) = (24) = 24\mathbb{Z}.
\]
The canonical homomorphism is:
\[
\phi \colon \faktor{\mathbb{Z}}{24\mathbb{Z}} \longrightarrow \faktor{\mathbb{Z}}{4\mathbb{Z}} \times \faktor{\mathbb{Z}}{6\mathbb{Z}}, \quad z + 24\mathbb{Z} \longmapsto (z + 4\mathbb{Z}, \, z + 6\mathbb{Z}).
\]
\begin{itemize}
    \item \textbf{Failure of Surjectivity:} Consider the element $(0 + 4\mathbb{Z}, 1 + 6\mathbb{Z}) \in \faktor{\mathbb{Z}}{4\mathbb{Z}} \times \faktor{\mathbb{Z}}{6\mathbb{Z}}$. If there existed a pre-image $z + 24\mathbb{Z}$, it would have to satisfy the system of congruences:
    \begin{align*}
    z &\equiv 0 \pmod 4 \implies z \text{ is even}, \\
    z &\equiv 1 \pmod 6 \implies z \text{ is odd},
    \end{align*}
    which is impossible. Hence:
    \[
    (0, 1) \notin \operatorname{Im}(\phi),
    \]
    so $\phi$ is not surjective.
    \item \textbf{Failure of Injectivity:} The element $12 + 24\mathbb{Z} \in \faktor{\mathbb{Z}}{24\mathbb{Z}}$ is non-zero since $24 \nmid 12$. However:
    \[
    \phi(12 + 24\mathbb{Z}) = (12 + 4\mathbb{Z}, \, 12 + 6\mathbb{Z}) = (0 + 4\mathbb{Z}, \, 0 + 6\mathbb{Z}) = (\bar{0}, \bar{0}).
    \]
    Thus, $12 + 24\mathbb{Z} \in \ker(\phi) \setminus \{0\}$, so $\phi$ is not injective.
\end{itemize}
\end{example}

\begin{steps}[Algebraic Reason for the Failure]
Notice that $| \faktor{\mathbb{Z}}{24\mathbb{Z}} | = 24$, whereas $| \faktor{\mathbb{Z}}{4\mathbb{Z}} \times \faktor{\mathbb{Z}}{6\mathbb{Z}} | = 4 \times 6 = 24$. Even though the two finite rings have the exact same cardinality, $\phi$ fails to be bijective because:
\begin{enumerate}[label=\textbf{(\alph*)}]
    \setcounter{enumi}{0} \item The sum of the ideals is $(4) + (6) = \gcd(4, 6)\mathbb{Z} = 2\mathbb{Z} \neq \mathbb{Z}$, so $4$ and $6$ are not coprime.
    \setcounter{enumi}{1} \item The kernel of the map $\phi$ is $I \cap J = \operatorname{lcm}(4, 6)\mathbb{Z} = 12\mathbb{Z}$, which strictly contains the product ideal $IJ = 24\mathbb{Z}$.
\end{enumerate}
\end{steps}
\end{content}

\subsection{Statement and Proof of the Chinese Remainder Theorem}

\begin{speech}[00:05:25 - 00:06:11]
Sorry. Maybe I don't want that shadow here. \inlinemetanote{adjusts the blackboard position}

Okay, the Chinese Remainder Theorem says: if we have the additional assumption that $I$ and $J$ are what we called coprime, that their sum is equal to the unit ideal:
\[
I + J = (1).
\]
\inlinemetanote{writes: CRT: If $I+J=(1)$} If this is true, if that holds, then $\phi$ is an isomorphism \inlinemetanote{writes: then $\phi$ is an isom.}.

So, in particular, it means it's both injective and surjective. And we have to prove both; we get neither of them for free, because, you know, they failed here. So, okay, the proof. \inlinemetanote{writes: Proof}

Let us start with $\phi$ is injective \inlinemetanote{writes: $\phi$ is injective}. So, we have to prove that.
\end{speech}

\begin{content}
\begin{theorem}[Chinese Remainder Theorem for Two Ideals]\label[theorem]{thm:crt-two-ideals}
Let $R$ be a commutative ring with identity, and let $I, J \subseteq R$ be ideals such that they are \newterm{comaximal} (or coprime), meaning:
\begin{equation}\label{eq:crt-comaximal-hypothesis}
I + J = (1) = R.
\end{equation}
Then the canonical ring homomorphism:
\begin{equation}\label{eq:crt-isomorphism}
\phi \colon \faktor{R}{IJ} \xrightarrow{\quad \sim \quad} \faktor{R}{I} \times \faktor{R}{J}, \quad \bar{z} \longmapsto (\bar{z}, \bar{z})
\end{equation}
is an isomorphism of rings.
\end{theorem}
\end{content}

\begin{hidden}
% timestamp: 00:06:15
% topic: Chinese Remainder Theorem for two ideals
% board_state: thm:crt-two-ideals, ex:non-coprime-crt
% next_goal: Prove injectivity of phi using the comaximal identity i + j = 1
% open_loops: none
\end{hidden}

\begin{proofWrapper}[Proof of \cref{thm:crt-two-ideals}]

\subsubsection{Injectivity of \texorpdfstring{$\phi$}{phi}}

\begin{speech}[00:06:11 - 00:07:50]
So, suppose we have $\bar{z}$ in $R \bmod IJ$ \inlinemetanote{writes: Suppose $\bar{z} \in \faktor{R}{IJ}$} such that $\phi(\bar{z}) = (\bar{0}, \bar{0})$ \inlinemetanote{writes: such that $\phi(\bar{z}) = (\bar{0}, \bar{0})$}. So, we just have to study the kernel for injectivity, okay?

Then what do we know? Since that means that $\bar{z} = 0 \bmod I$ and $z = 0 \bmod J$ \inlinemetanote{writes: Since $\bar{z} = 0 \bmod I$ and $z = 0 \bmod J$}. That's what those two zeros mean, okay? You agree?

And so, this means that $z$ is in $I$ and $z$ is in $J$ \inlinemetanote{writes \qt{$z \in I$} and \qt{$z \in J$}}. That's how that's how you're zero $\bmod I$ and $J$.

Okay, but we need to show that $z$ is in $IJ$. We have to somehow use the fact that $z$ is in $I$ and $z$ is in $J$ to show it's in $IJ$, which is a hard thing to do in general, not true; but we have to use this equation. So, but we know --- we haven't used it yet, but we know --- that there exists some $i \in I$ and there exists a $j \in J$ such that $i + j = 1$ \inlinemetanote{writes: But we know $\exists i \in I, j \in J$ such that $i+j=1$}. That's the kind of magic of this assumption. That's exactly what this assumption says: there's an $i$ in $I$, a little $i$ in $I$, and a little $j$ in $J$ such that they sum to $1$.
\end{speech}

\begin{speech}[00:07:50 - 00:09:25]
And the whole idea of this injectivity is just one idea. Up to here there have been no ideas. And you just multiply this by $z$. So, the idea is: multiply this by $z$ \inlinemetanote{writes: $\cdot z$}. And $iz + jz$ --- it's a very simple idea --- is equal to $z$ \inlinemetanote{writes: $iz + jz = z$}. There.

So, that's undeniable: I multiply by $z$, I get this. And then I look at this equation: this is $z$... This one, this $i$ is in $I$, but $z$ we know is already in $J$ \inlinemetanote{draws arrows showing $i \in I$ and $z \in J$}. This $j$ is in $J$, but $z$ we know is in $I$ \inlinemetanote{draws arrows showing $j \in J$ and $z \in I$}. So, both of these things on the right are in $IJ$, in different ways. We use the fact that $z$ is both in $I$ and $J$ in this analysis.

And if I look at it, that means this term is in $IJ$, that term is in $IJ$, and since $IJ$ is an ideal, the conclusion is $z \in IJ$ \inlinemetanote{writes: $z \in IJ$}. That's the conclusion, which is the one we wanted, because if $z$ is in $IJ$, then it was zero up here. And this implies, so, $\bar{z} = 0$ in $R \bmod IJ$ \inlinemetanote{writes: So $\bar{z} = 0 \in \faktor{R}{IJ}$}.

So, that shows it's injective. And as I said, this whole proof only has one idea. Here, just following your nose, writing the data, and the only idea is: you take this, multiply by $z$, that's it.
\end{speech}

\begin{content}[Proof of Injectivity]
\begin{proof}[Verification of Injectivity]
To prove that $\phi$ is injective, we study its kernel. Suppose $\bar{z} = z + IJ \in \faktor{R}{IJ}$ satisfies:
\[
\phi(\bar{z}) = (\bar{0}, \bar{0}) \in \faktor{R}{I} \times \faktor{R}{J}.
\]
By the definition of the component classes, this means:
\[
z \equiv 0 \pmod I \implies z \in I, \qquad \text{and} \qquad z \equiv 0 \pmod J \implies z \in J.
\]
Thus, $z \in I \cap J$. To show that $\bar{z} = 0 \in \faktor{R}{IJ}$, we must prove that $z \in IJ$.

\par\vspace{1ex}\noindent
By assumption, $I + J = (1)$, so there exist elements $i \in I$ and $j \in J$ such that:
\begin{equation}\label{eq:crt-unit-decomp}
i + j = 1.
\end{equation}
Multiplying both sides of \eqref{eq:crt-unit-decomp} by $z$ yields:
\begin{equation}\label{eq:crt-injectivity-identity}
z = (i + j)z = iz + jz.
\end{equation}
We now analyze each of the two summands in \eqref{eq:crt-injectivity-identity}:
\begin{itemize}
    \item \textbf{First Summand ($iz$):} In the product $iz$, we have $i \in I$ and $z \in J$, so $iz \in IJ$.
    \item \textbf{Second Summand ($jz$):} In the product $jz$, we have $j \in J$ and $z \in I$, so $jz \in JI = IJ$ (by commutativity of $R$).
\end{itemize}
Since $IJ$ is an ideal, it is closed under addition, which implies:
\[
z = \underbrace{iz}_{\in IJ} + \underbrace{jz}_{\in IJ} \in IJ.
\]
Therefore, $z \in IJ$, which implies $\bar{z} = 0$ in $\faktor{R}{IJ}$. This shows that $\ker(\phi) = \{0\}$, so $\phi$ is injective.
\end{proof}
\end{content}

\subsubsection{Surjectivity of \texorpdfstring{$\phi$}{phi}}

\begin{speech}[00:09:25 - 00:11:15]
Okay, but as I said, we don't get anything for free. Now we have to prove surjectivity, which looks a little more complicated.

So, now we have to prove that this map is surjective \inlinemetanote{writes: $\phi$ is surjective}.

Surjective means we start: so, choose any $\bar{x} \in R \bmod I$ and $\bar{y} \in R \bmod J$ \inlinemetanote{writes: Choose any $\bar{x} \in \faktor{R}{I}, \bar{y} \in \faktor{R}{J}$}, right? Then we have to find... we have to construct a $\bar{z} \in R \bmod IJ$ such that \inlinemetanote{writes: We have to construct a $\bar{z} \in \faktor{R}{IJ}$ such that} $z$ is congruent to $x \bmod I$ and $z$ is congruent to $y \bmod J$ \inlinemetanote{writes \qt{$z \equiv x \bmod I$} and \qt{$z \equiv y \bmod J$}}. If I do that, that's surjectivity. That's the Chinese Remainder problem that we talked about at the beginning. So, you pick where that you want, and you have to find the $z$ that satisfies that.

Okay, so, how are we going to do that? Well, what I know is $i + j = 1$ \inlinemetanote{writes: $i+j=1$}. That's all I know. So, now I need to do something. And what do you suggest I do here?

Well, before, the only thing we had was $z$. That was our putative element of the kernel. And we took this equation, we multiplied by $z$. Now, what we have in our hands are $x$ and $y$. So, we have this equation, so it's more or less reasonable we multiply this by $x$ and $y$. What else could we do? Let's multiply it by $x$ \inlinemetanote{writes: $\cdot x$}.
\end{speech}

\begin{speech}[00:11:15 - 00:12:57]
$ix + jx = x$ \inlinemetanote{writes: $ix + jx = x$}. That's an interesting equation.

If I look at this term here, $jx$ \inlinemetanote{underlines $jx$}, this is congruent to $x \bmod I$ \inlinemetanote{writes: $jx \equiv x \bmod I$}, because this equation says that if I look $\bmod I$, this is congruent... it's to $x$, okay? So, just this equation shows that $jx$ is congruent to $x \bmod I$. Is that clear?

And $jx$, that's congruent to $0 \bmod J$ \inlinemetanote{writes: $jx \equiv 0 \bmod J$}, because it has a $j$ in it. It's in the ideal.

All right, so, that was some progress. You can imagine now we should multiply it by $y$ \inlinemetanote{writes \qt{$i+j=1$} and \qt{$\cdot y$}}: $iy + jy = y$ \inlinemetanote{writes: $iy + jy = y$}. And now we want to look at this one \inlinemetanote{boxes $iy$}: $iy$ is congruent to $0 \bmod I$ \inlinemetanote{writes: $iy \equiv 0 \bmod I$}, because it has an $i$ in it. But $iy$ is congruent to $y \bmod J$ \inlinemetanote{writes: $iy \equiv y \bmod J$}, because of this equation. If I look at this $\bmod J$, that goes away, and $iy$ is congruent...

So, that's perfect. We've used this... there was only one trick in this whole thing. We used it once for injectivity, and we're going to use it twice for surjectivity. So, what should $z$ be here? Then I claim that we're now prepared to construct $z$. What should I write $z$ here? Yeah?
\end{speech}

\begin{student}[Student Answer]
$jx + iy$.
\end{student}

\begin{speech}[00:12:59 - 00:13:45]
Yeah, did you get it right? $jx + iy$, yes: $jx + iy$ \inlinemetanote{writes: $\implies z = jx + iy$}. Because this thing, you just add them up now. It's $x \bmod I$ and $y \bmod J$, that's it.

So, that's the end of surjectivity. So, here for this one, $\phi(\bar{z})$ is exactly equal to $(\bar{x}, \bar{y})$ \inlinemetanote{writes: $\phi(\bar{z}) = (\bar{x}, \bar{y})$}. And that's the end of the proof of the Chinese Remainder Theorem. And the proof... so, it's not a proof with no ideas. In some sense, it's one idea used three times.

But if you, I mean, if you understand this proof, you'll never forget it for the rest of your life, I think, somehow, maybe, I don't know.
\end{speech}

\begin{content}[Proof of Surjectivity and Full Restatement]
\begin{proof}[Verification of Surjectivity]
Let $\bar{x} = x + I \in \faktor{R}{I}$ and $\bar{y} = y + J \in \faktor{R}{J}$ be arbitrary elements. We must construct an element $z \in R$ satisfying the simultaneous congruences:
\begin{equation}\label{eq:crt-congruences-goal}
z \equiv x \pmod I \qquad \text{and} \qquad z \equiv y \pmod J.
\end{equation}
We use the elements $i \in I$ and $j \in J$ with $i + j = 1$:
\begin{itemize}
    \item \textbf{Action on $x$:} Multiplying $i + j = 1$ by $x$ gives $ix + jx = x$. Analyzing the term $jx$:
    \begin{align}
    jx &= x - ix \equiv x \pmod I && (\text{since } i \in I \implies ix \in I), \label{eq:jx-mod-I} \\
    jx &\equiv 0 \pmod J && (\text{since } j \in J \implies jx \in J). \label{eq:jx-mod-J}
    \end{align}
    \item \textbf{Action on $y$:} Multiplying $i + j = 1$ by $y$ gives $iy + jy = y$. Analyzing the term $iy$:
    \begin{align}
    iy &\equiv 0 \pmod I && (\text{since } i \in I \implies iy \in I), \label{eq:iy-mod-I} \\
    iy &= y - jy \equiv y \pmod J && (\text{since } j \in J \implies jy \in J). \label{eq:iy-mod-J}
    \end{align}
\end{itemize}
Now define:
\begin{equation}\label{eq:crt-explicit-z}
z := jx + iy.
\end{equation}
Evaluating $z$ modulo each ideal:
\begin{align*}
z &\equiv jx + iy \equiv x + 0 \equiv x \pmod I, \\
z &\equiv jx + iy \equiv 0 + y \equiv y \pmod J.
\end{align*}
Thus, $\phi(z + IJ) = (z + I, z + J) = (x + I, y + J) = (\bar{x}, \bar{y})$. This proves that $\phi$ is surjective.
Having established both injectivity and surjectivity, $\phi$ is an isomorphism of rings.
\end{proof}
\end{content}

\end{proofWrapper}

\begin{insight}[The Partition of Unity Principle in the Chinese Remainder Theorem]
The comaximal equation $i + j = 1$ with $i \in I$ and $j \in J$ functions as an algebraic \qt{partition of unity}:
\begin{itemize}
    \item \textbf{Indicator for $I$:} The element $j$ acts as an indicator for the ideal $I$: it satisfies $j \equiv 1 \pmod I$ and $j \equiv 0 \pmod J$.
    \item \textbf{Indicator for $J$:} The element $i$ acts as an indicator for the ideal $J$: it satisfies $i \equiv 0 \pmod I$ and $i \equiv 1 \pmod J$.
\end{itemize}
The explicit solution $z = jx + iy$ linearly combines these indicator elements with the target values $x$ and $y$. This algebraic construction mirrors orthogonal projection bases in linear algebra and Lagrange interpolation polynomials in numerical analysis.
\end{insight}

\begin{hidden}
% timestamp: 00:13:45
% topic: Completed proof of Chinese Remainder Theorem
% board_state: thm:crt-two-ideals, eq:crt-explicit-z
% next_goal: Applications to the integers Z and prime power decomposition
% open_loops: none
\end{hidden}

\subsection{The Chinese Remainder Theorem in \texorpdfstring{$\mathbb{Z}$}{Z}}

\begin{speech}[00:13:45 - 00:15:00]
And you can... yeah. This is how it works, the Chinese Remainder Theorem. You just take that $i + j = 1$ and use it. \inlinemetanote{erases the lower board}

Okay, so, let's go back to the integers, because you can use this now many, many times, and you get some fun results about the integers, which, you know, are very closely related to topics in school, but maybe usually not made explicit there.

So, this says that if I just write this theorem, the Chinese Remainder Theorem for $R = \mathbb{Z}$ \inlinemetanote{writes: CRT in $R = \mathbb{Z}$}, then it says that if $a$ and $b$ are relatively prime \inlinemetanote{writes: $(a, b) \text{ co-pr.}$} --- this is in the language of the Chinese Remainder Theorem, this says that the ideal of $a$ plus the ideal of $b$ is equal to $(1)$ \inlinemetanote{writes: $( (a) + (b) = (1) )$} ---, then $\mathbb{Z} \bmod ab$ is isomorphic to $\mathbb{Z} \bmod a \times \mathbb{Z} \bmod b$ \inlinemetanote{writes: then $\faktor{\mathbb{Z}}{(ab)} \cong \faktor{\mathbb{Z}}{(a)} \times \faktor{\mathbb{Z}}{(b)}$}, okay?
\end{speech}

\begin{content}
\begin{corollary}[Chinese Remainder Theorem for Coprime Integers]\label[corollary]{cor:crt-integers}
Let $R = \mathbb{Z}$, and let $a, b \in \mathbb{Z}$ be coprime integers, i.e., $\gcd(a, b) = 1$. By Bézout's identity, this is equivalent to comaximal principal ideals:
\[
(a) + (b) = (\gcd(a, b)) = (1) = \mathbb{Z}.
\]
Since the product of principal ideals in $\mathbb{Z}$ is $(a)(b) = (ab)$, \cref{thm:crt-two-ideals} yields a ring isomorphism:
\begin{equation}\label{eq:crt-integers-formula}
\faktor{\mathbb{Z}}{ab\mathbb{Z}} \xrightarrow{\quad \sim \quad} \faktor{\mathbb{Z}}{a\mathbb{Z}} \times \faktor{\mathbb{Z}}{b\mathbb{Z}}, \quad z \bmod ab \longmapsto (z \bmod a, \, z \bmod b).
\end{equation}
\end{corollary}
\end{content}

\begin{speech}[00:15:00 - 00:17:15]
So, you think, okay, but this is only two, what if I factor in more? But you can do them one at a time. So, for example, if I take some number like $3 \times 5 \times 7$, that has... it has $105$, it has three factors \inlinemetanote{writes: $3 \cdot 5 \cdot 7 = 105$}.

So, how should I do with $a$, $b$, and $c$? Well, I can always group it as $a = 3$, $b = 35$ \inlinemetanote{writes: $a = 3, b = 35$}. And then the first step is that $\mathbb{Z} \bmod 105$ is isomorphic to $\mathbb{Z} \bmod 3 \times \mathbb{Z} \bmod 35$ \inlinemetanote{writes: $\faktor{\mathbb{Z}}{(105)} \cong \faktor{\mathbb{Z}}{(3)} \times \faktor{\mathbb{Z}}{(35)}$}, right?

And then I can again now do it to the $35$ and split it further. So, this further isomorphic, I just apply it now to the $35$ \inlinemetanote{writes: $\cong \faktor{\mathbb{Z}}{(3)} \times \faktor{\mathbb{Z}}{(5)} \times \faktor{\mathbb{Z}}{(7)}$}. We could have stated this at the very beginning, but there's no reason to. So, it says I can just write them all out as long as I'm at every time splitting coprime factors out. Is that clear?

And so, then the fanciest way to write it is: suppose I have $m$ that's an integer... let's call it $n$. There's a prime factorization: $p_1^{k_1} \dots$, where these are the distinct primes that come here, $p_r^{k_r}$ \inlinemetanote{writes: $n = p_1^{k_1} p_2^{k_2} \cdots p_r^{k_r}$}. So, this is the prime factorization of $n$. It has some distinct primes to various powers. There's a unique way to write that. That's the unique factorization of primes.

And then if I just apply this, I'm not allowed to separate the prime factors, because if I separate $p$ and, you know, $p^2$ and $p$, they're not relatively prime. But I can separate different prime factors. So, what this shows, this Chinese Remainder Theorem shows that this is isomorphic to the product... \inlinemetanote{writes: $\faktor{\mathbb{Z}}{(n)} \cong \faktor{\mathbb{Z}}{(p_1^{k_1})} \times \dots \times \faktor{\mathbb{Z}}{(p_r^{k_r})}$}. I just separate out the primes. The primes all come apart.

Okay? So, I don't know whether this was... I mean, whatever happens in school is usually not in this language.
\end{speech}

\begin{content}[Iterated Factoring and Prime-Power Decomposition]
\begin{example}[Iterated Decomposition of \texorpdfstring{$\faktor{\mathbb{Z}}{(105)}$}{Z/(105)}]\label[example]{ex:crt-105}
Consider $n = 105 = 3 \cdot 5 \cdot 7$. Grouping the factors as $a = 3$ and $b = 35$, with $\gcd(3, 35) = 1$:
\[
\faktor{\mathbb{Z}}{105\mathbb{Z}} \cong \faktor{\mathbb{Z}}{3\mathbb{Z}} \times \faktor{\mathbb{Z}}{35\mathbb{Z}}.
\]
Since $\gcd(5, 7) = 1$, we decompose $\faktor{\mathbb{Z}}{35\mathbb{Z}} \cong \faktor{\mathbb{Z}}{5\mathbb{Z}} \times \faktor{\mathbb{Z}}{7\mathbb{Z}}$, which yields:
\[
\faktor{\mathbb{Z}}{105\mathbb{Z}} \cong \faktor{\mathbb{Z}}{3\mathbb{Z}} \times \faktor{\mathbb{Z}}{5\mathbb{Z}} \times \faktor{\mathbb{Z}}{7\mathbb{Z}}.
\]
\end{example}

\begin{corollary}[Prime-Power Factorization of \texorpdfstring{$\faktor{\mathbb{Z}}{n\mathbb{Z}}$}{Z/nZ}]\label[corollary]{cor:crt-prime-power-decomp}
Let $n \in \mathbb{N}_{\ge 2}$ have the prime factorization into distinct primes $p_1 < p_2 < \dots < p_r$:
\begin{equation}\label{eq:n-prime-factorization}
n = p_1^{k_1} p_2^{k_2} \cdots p_r^{k_r}, \quad k_i \ge 1.
\end{equation}
Since $\gcd(p_i^{k_i}, p_j^{k_j}) = 1$ for all $i \neq j$, induction on $r$ gives the canonical ring isomorphism:
\begin{equation}\label{eq:crt-prime-power-isom}
\faktor{\mathbb{Z}}{n\mathbb{Z}} \xrightarrow{\quad \sim \quad} \faktor{\mathbb{Z}}{p_1^{k_1}\mathbb{Z}} \times \faktor{\mathbb{Z}}{p_2^{k_2}\mathbb{Z}} \times \dots \times \faktor{\mathbb{Z}}{p_r^{k_r}\mathbb{Z}} = \prod_{i=1}^r \faktor{\mathbb{Z}}{p_i^{k_i}\mathbb{Z}}.
\end{equation}
\end{corollary}

\begin{steps}[Why Prime Powers Cannot Be Split Further]
While distinct prime powers $p_i^{k_i}$ and $p_j^{k_j}$ are coprime and can be separated, factors involving the same prime cannot be split via the Chinese Remainder Theorem. For example:
\[
\faktor{\mathbb{Z}}{p^2\mathbb{Z}} \not\cong \faktor{\mathbb{Z}}{p\mathbb{Z}} \times \faktor{\mathbb{Z}}{p\mathbb{Z}},
\]
because $\gcd(p, p) = p \neq 1$. The ring $\faktor{\mathbb{Z}}{p^2\mathbb{Z}}$ contains non-zero nilpotent elements ($p + p^2\mathbb{Z} \neq 0$, but $(p + p^2\mathbb{Z})^2 = 0$), whereas $\faktor{\mathbb{Z}}{p\mathbb{Z}} \times \faktor{\mathbb{Z}}{p\mathbb{Z}}$ is a direct product of fields and has no non-zero nilpotents.
\end{steps}
\end{content}

\begin{hidden}
% timestamp: 00:18:00
% topic: Units of Z/nZ and Euler's totient function
% board_state: cor:crt-prime-power-decomp, eq:n-prime-factorization
% next_goal: Deduce Euler's totient formula phi(n) and discuss Bernoulli number B_12
% open_loops: none
\end{hidden}

\subsection{Units in \texorpdfstring{$\faktor{\mathbb{Z}}{n\mathbb{Z}}$}{Z/nZ} and Euler's Totient Function}

\begin{speech}[00:17:15 - 00:19:16]
Now we can consider the units. So, we have this idea of the units. We consider the units of this ring \inlinemetanote{writes: $\operatorname{Unit}(\faktor{\mathbb{Z}}{(n)})$} $\mathbb{Z} \bmod n\mathbb{Z}$... Let me just write it, because I don't want to confuse the notation: the units of this ring $\mathbb{Z} \bmod n\mathbb{Z}$, $\mathbb{Z} \bmod n$.

The units. Those are the invertible elements here. But remember, $1$ here is just $(1, 1)$, everybody is just $1$, the unit here is the unit everywhere. So, how can you be invertible on this side? Well, you have to be invertible in every factor. So, this says the units here are also just the product here: it's the product of the units \inlinemetanote{writes: $= \operatorname{Unit}(\faktor{\mathbb{Z}}{(p_1^{k_1})}) \times \dots \times \operatorname{Unit}(\faktor{\mathbb{Z}}{(p_r^{k_r})})$}.

\inlinemetanote{writes \qt{$\operatorname{Unit}(\faktor{\mathbb{Z}}{(p_r^{k_r})})$} on the far right} This was bad board management.

$p_r^{k_r}$, okay? So, if you're interested in the units of this ring, this is extremely useful, because it just tells you the units are just the product. You have to pick a unit here, a unit here, a unit here, that's all it is. And moreover, then it also tells you how many there are!

Because I hinted at this last... whenever we discussed this a few weeks ago. But this is related to this Euler totient function. That Euler totient function is the number of units. And this says if you want to know the number of units of $\mathbb{Z} \bmod n\mathbb{Z}$, it's just the product of the number of units of all of these prime ones, prime power ones. But it's easy to calculate the number of units in the prime power, which I just remind you, I think... maybe I left it as an exercise. If so, I'll leave it as an exercise again.
\end{speech}

\begin{content}[Group of Units of the Direct Product]
\begin{proposition}[Multiplicative Group of Units of a Product Ring]\label[proposition]{prop:unit-group-product}
Let $R_1, \dots, R_r$ be commutative rings with identity. An element $(u_1, \dots, u_r) \in \prod_{i=1}^r R_i$ is a unit if and only if each component $u_i \in \operatorname{Unit}(R_i)$ is a unit. Therefore, taking unit groups commutes with finite direct products:
\begin{equation}\label{eq:unit-group-isomorphism}
\operatorname{Unit}\left(\prod_{i=1}^r R_i\right) \cong \prod_{i=1}^r \operatorname{Unit}(R_i).
\end{equation}
Applying this to the prime-power decomposition \eqref{eq:crt-prime-power-isom} of $\faktor{\mathbb{Z}}{n\mathbb{Z}}$:
\begin{equation}\label{eq:units-zn-formula}
\operatorname{Unit}\left(\faktor{\mathbb{Z}}{n\mathbb{Z}}\right) \cong \operatorname{Unit}\left(\faktor{\mathbb{Z}}{p_1^{k_1}\mathbb{Z}}\right) \times \operatorname{Unit}\left(\faktor{\mathbb{Z}}{p_2^{k_2}\mathbb{Z}}\right) \times \dots \times \operatorname{Unit}\left(\faktor{\mathbb{Z}}{p_r^{k_r}\mathbb{Z}}\right).
\end{equation}
\end{proposition}
\end{content}

\begin{speech}[00:19:16 - 00:20:45]
If I take $\phi$ of a power of a prime \inlinemetanote{erases part of the lower board and writes: $\phi(p^k) =$} --- and let me just follow the notation here: $p$, but any prime, no index, with a $k$ ---, so, what is this?

Well, this is all the numbers in $\mathbb{Z} \bmod p^k \mathbb{Z}$ that are relatively... all the numbers that are relatively prime to $p^k$ that are less than $p^k$. And how are you going to be a common factor? You have to have a $p$ in it. And so, how many people have $p$'s in it? Well, exactly $1/p$ of these has $p$ in it. So, the answer is $p^k - p^{k-1}$ \inlinemetanote{writes: $= p^k - p^{k-1}$}, or you can write this as $p^{k-1}(p - 1)$ \inlinemetanote{writes: $= p^{k-1}(p - 1)$}. Okay?

So, that gives a solution, a kind of formula to this Euler totient. It says that if I want to know what this totient function is, that's the order of the units \inlinemetanote{writes: $\phi(n) = |\operatorname{Unit}(\faktor{\mathbb{Z}}{n})| =$}, this is equal to the product from $i = 1$ to $r$, in that language, of $p_i^{k_i - 1}(p_i - 1)$ \inlinemetanote{writes: $= \prod_{i=1}^r p_i^{k_i - 1}(p_i - 1)$}. That's a pretty good formula for it, probably the fastest way to calculate it.

Uh, but it's a confusing formula, like you might think... okay, so, this is the part that I view as perhaps related to things that may have been in school.
\end{speech}

\recap{Euler's totient function was introduced in Week~3 (\cref{def:eulers-totient}), with its values $\phi(p)$ and $\phi(p^2)$ (\cref{prop:phi-prime-powers}); the units of $\faktor{\mathbb{Z}}{m\mathbb{Z}}$ were determined in Week~2 (\cref{thm:units-z-mod-m}). The Chinese Remainder Theorem now gives the general formula.}

\begin{content}[Derivation of Euler's Product Formula]
\begin{definition}[Euler's Totient Function]\label[definition]{def:euler-totient}
The \newterm{Euler totient function} $\varphi(n)$ (written on the blackboard as $\phi(n)$) counts the number of invertible elements in the ring $\faktor{\mathbb{Z}}{n\mathbb{Z}}$:
\begin{equation}\label{eq:totient-def-w5wed}
\varphi(n) := \left| \operatorname{Unit}\left(\faktor{\mathbb{Z}}{n\mathbb{Z}}\right) \right| = \left| \{ m \in \{1, \dots, n\} \mid \gcd(m, n) = 1 \} \right|.
\end{equation}
\end{definition}

\begin{theorem}[Euler's Product Formula]\label[theorem]{thm:euler-totient-product}
For a prime power $p^k$ ($k \ge 1$), the elements in $\{1, 2, \dots, p^k\}$ that are \emph{not} coprime to $p^k$ are precisely the multiples of $p$:
\[
\{ p, \, 2p, \, 3p, \, \dots, \, p^{k-1} p = p^k \}.
\]
There are exactly $\frac{p^k}{p} = p^{k-1}$ such multiples. Subtracting these from the total of $p^k$ elements gives:
\begin{equation}\label{eq:phi-prime-power}
\varphi(p^k) = p^k - p^{k-1} = p^{k-1}(p - 1) = p^k \left(1 - \frac{1}{p}\right).
\end{equation}
Since the order of a direct product of finite groups is the product of their orders, taking cardinalities in \eqref{eq:units-zn-formula} yields:
\begin{equation}\label{eq:euler-general-formula}
\varphi(n) = \prod_{i=1}^r \varphi\left(p_i^{k_i}\right) = \prod_{i=1}^r p_i^{k_i - 1}(p_i - 1) = n \prod_{i=1}^r \left(1 - \frac{1}{p_i}\right).
\end{equation}
\end{theorem}
\end{content}

\begin{speech}[00:20:45 - 00:22:20]
It's a formula that almost looks like if you know what this totient thing, since all the primes are here, it seems like you could reconstruct the actual number, because all its prime factorization is here. But it's not true! It's because of these $(p_i - 1)$'s. They confuse the issue. Like if you take a prime number... well, I can tell you a personal thing: my favorite prime number.

If I take $691$ \inlinemetanote{writes: $\phi(691) =$} --- which I assume you know is a prime number; this is an extremely famous prime number in mathematics ---, but this is a prime number, and if I take that, it's just the $p - 1$, so, this is $690$ \inlinemetanote{writes: $= 690$}, and that has a lot of factors in it. It has, you know, some $3$ and $23$, $5$ and $2$ \inlinemetanote{writes: $= 3 \cdot 23 \cdot 5 \cdot 2$}. So, you somehow lose... If you start with some nice prime number, this totient function gives you lots of primes and they get mixed up here. So, you can't recover, even though it almost looks like it, you can't recover the $n$ from the totient value, because this $p - 1$ kind of factors everything out.

Have you ever heard of this prime number? It's like one of the most famous prime numbers in mathematics. It comes from Bernoulli $12$ \inlinemetanote{writes: $B_{12} =$} which is just a great Bernoulli number. It's $-691 / 2730$ \inlinemetanote{writes: $-\frac{691}{2730}$}. But have you learned about the Bernoulli numbers? They're also a local product, made in Basel. Because the Bernoullis were in Basel, that's why.

Okay, so, I think this is more or less what I wanted to say about the Chinese Remainder Theorem.
\end{speech}

\begin{content}[The Irregular Prime 691 and Bernoulli Numbers]
\begin{example}[Euler's Totient on the Prime 691]\label[example]{ex:prime-691}
For the prime $p = 691$, the value of Euler's totient function is:
\begin{equation}\label{eq:phi-691}
\varphi(691) = 691 - 1 = 690 = 2 \cdot 3 \cdot 5 \cdot 23.
\end{equation}
The prime $691$ is famous in algebraic number theory as the numerator of the 12th Bernoulli number $B_{12}$:
\begin{equation}\label{eq:bernoulli-12}
B_{12} = -\frac{691}{2730}.
\end{equation}
\end{example}

\begin{steps}[Irreversibility of Euler's Totient Function]
Even though the prime factorization of $n$ determines $\varphi(n)$ via formula \eqref{eq:euler-general-formula}, one cannot generally reconstruct $n$ from $\varphi(n)$. The presence of the terms $(p_i - 1)$ mixes new prime factors (such as $23$ appearing in $\varphi(691) = 690$) that have no relation to the original prime $p = 691$. In modern cryptography (e.g., the RSA cryptosystem), the computational difficulty of finding $\varphi(n) = (p - 1)(q - 1)$ without knowing the prime factors $p$ and $q$ of a large integer $n = pq$ provides the trapdoor mechanism for secure encryption.
\end{steps}
\end{content}

\begin{aiNote}[Beyond the Lecture: Kummer's Regularity Criterion and $B_{12}$]
% [PERSISTENT, ADDITION]
The prime $691$ is the prototypical example of an \newterm{irregular prime}. In 1850, Ernst Kummer proved Fermat's Last Theorem for all regular prime exponents. Kummer showed that a prime $p$ is regular if and only if $p$ does not divide the numerator of any of the even Bernoulli numbers $B_2, B_4, \dots, B_{p-3}$. Since $691$ divides the numerator of $B_{12}$, it is irregular, and it plays a prominent role in modern modular forms and arithmetic geometry (e.g., Ramanujan's $\tau$-function congruences $\tau(n) \equiv \sigma_{11}(n) \pmod{691}$).
\end{aiNote}

\section{Polynomial Rings over Fields}

\subsection{Definitions and the Degree Function}

\begin{speech}[00:22:20 - 00:23:46]
And so, the last topic of the week, which is going to go on to next week, has to do with polynomials \inlinemetanote{writes: Polynomials}.

And we've talked about these things in various ways, but I'll try to say it now clearly from the beginning. So, we start with $K$ is a field \inlinemetanote{writes: $K$ is a field}.

And we've discussed before there is a ring which is polynomials in one variable \inlinemetanote{writes: $K[x]$ polys in a single variable} with coefficients in $K$ \inlinemetanote{writes: with coeffs in $K$}.

Okay. And now here is something that one has to just state a little bit carefully: it's the degree of a polynomial. We've said it kind of informally, but now I'm going to say it carefully, and there's only one point that's... All right, so, we would like to define the degree of a polynomial \inlinemetanote{writes: $\deg(f(x))$}. This seems kind of silly to even talk about, but we're going to say. So, we have some element of this ring \inlinemetanote{writes: $f(x) \in K[x]$}, and we want to find the degree of the polynomial.
\end{speech}

\begin{hidden}
% timestamp: 00:05:00
% topic: Definition of polynomial rings K[x] and the degree function
% board_state: def:polynomial-ring, def:degree-polynomial
% next_goal: Define the degree for non-zero polynomials and discuss deg(0)
% open_loops: none
\end{hidden}

\begin{speech}[00:23:46 - 00:25:16]
And so, the degree of the polynomial is that you write the polynomial, you find some coefficient, the highest coefficient which is not zero \inlinemetanote{writes \qt{$c_d x^d$} above}, so, this is the highest coefficient \inlinemetanote{writes: highest coeff.} which is not zero \inlinemetanote{writes: which is not zero}.

And that makes sense, because a polynomial is just some list of these coefficients. You find the highest one that's not zero, you look at the $d$ index, and that's the degree. So, this is the degree of this polynomial. Another way to write it informally: the degree of $f$ is equal to this $d$ \inlinemetanote{writes: $\deg(f) = d$}, where $d$ is the coefficient of the highest not zero. Okay, is that clear?

And the only subtle point of this is: what is the degree of the zero polynomial? \inlinemetanote{writes: $\deg(0)$} That is a super unpleasant thing to think about, because it has no coefficient that's not zero.
\end{speech}

\begin{speech}[00:25:16 - 00:26:46]
So, there's some kind of conventions about it, but the convention that's correct is that this is undefined \inlinemetanote{writes: undefined}.

You can make up your own conventions, but the convention that works the best is it's undefined. So, that was my whole point of having the degree, right? It's kind of an annoying thing, but you see that this recipe makes no sense, because you're supposed to write it out, find the highest coefficient which is not zero, and say that's the degree. But if it's zero, all the coefficients are zero, so, you're kind of, you know...

Anyway, so, it's an important thing and it sounds silly, and in a sense it's silly, but when you talk about degrees, you're not allowed to talk about the degree of the zero polynomial.

The degree of a constant is okay. The degree of a constant, like the degree of the polynomial $7$, as long as that $7$ is not zero, that has degree zero. But the degree of zero, you don't talk about.
\end{speech}

\begin{speech}[00:26:46 - 00:27:36]
I hope that wasn't too polemic, but that's just the standard mathematical definition, and since I'm supposed to be teaching you mathematics, I tell you the definition.

So, that should have been also when you learned about polynomials in school, you should have asked your teacher: what is the degree of the zero polynomial? You should have asked that question. \inlinemetanote{erases the middle chalkboard to the left}

I think that's the only thing I'm going to say about polynomials that most people don't learn in school. The rest I'm going to talk about today is stuff you've learned in school, but this particular point is that...
\end{speech}

\recap{Polynomial rings were defined in Week~3 (\cref{def:poly-ring-informal}, \cref{def:polynomial-ring-w3wed}), division with remainder in $\mathbb{Q}[x]$ was used in Week~4 (\cref{lem:poly-division-algorithm}), and division with remainder for polynomials is among the prerequisites (\cref{fact:pre-poly-division}). The lecturer now treats $K[x]$ over a field systematically.}

\begin{content}[Polynomial Ring and the Degree of a Polynomial]
\begin{definition}[Polynomial Ring over a Field]\label[definition]{def:poly-ring-k}
Let $K$ be a field. The \newterm{polynomial ring in one variable} over $K$, denoted by $K[x]$, is the set of formal expressions:
\begin{equation}\label{eq:poly-general-form}
f(x) = \sum_{i=0}^n c_i x^i = c_0 + c_1 x + c_2 x^2 + \dots + c_n x^n, \quad c_i \in K,
\end{equation}
equipped with standard polynomial addition and Cauchy product multiplication.
\end{definition}

\begin{definition}[Degree of a Polynomial]\label[definition]{def:poly-degree}
Let $f(x) \in K[x]$.
\begin{enumerate}[label=\textbf{(\arabic*)}]
    \setcounter{enumi}{0} \item \textbf{Non-Zero Polynomials:} If $f(x) \neq 0$, write $f(x) = \sum_{i=0}^n c_i x^i$. The \newterm{degree} of $f$, denoted by $\deg(f)$ or $\deg(f(x))$, is the highest index $d \in \mathbb{N}_0$ such that the coefficient $c_d \neq 0$:
    \begin{equation}\label{eq:deg-formula}
    \deg(f) := d \iff c_d \neq 0 \quad \text{and} \quad c_i = 0 \quad \forall i > d.
    \end{equation}
    The non-zero element $c_d \in K^\times$ is called the \newterm{leading coefficient}.
    \setcounter{enumi}{1} \item \textbf{Non-Zero Constants:} For any non-zero constant $c \in K^\times$, we have $\deg(c) = 0$.
    \setcounter{enumi}{2} \item \textbf{The Zero Polynomial:} For the zero polynomial $f(x) = 0$, the degree is explicitly:
    \begin{equation}\label{eq:deg-zero}
    \deg(0) := \text{undefined}.
    \end{equation}
\end{enumerate}
\end{definition}

\begin{remark}[Conventions for the Degree of the Zero Polynomial]\label[remark]{rem:deg-zero-convention}
In some textbook literature, the convention $\deg(0) := -\infty$ is adopted so that the degree formula $\deg(f \cdot g) = \deg(f) + \deg(g)$ remains formally valid when one factor is zero (with $-\infty + d = -\infty$). However, following the lecturer's strict convention, the degree of the zero polynomial is left strictly \textbf{undefined}.
\end{remark}

\begin{steps}[Why the Leading Coefficient Recipe Fails for Zero]
The algorithmic recipe for the degree requires identifying the \qt{highest non-zero coefficient $c_d$}. For the zero polynomial $0 \in K[x]$, all coefficients in $0 = \sum 0 \cdot x^i$ are identically zero, so no such index $d$ exists.
\end{steps}
\end{content}

\subsection{Polynomial Long Division}

\begin{speech}[00:27:36 - 00:29:16]
Okay, so, very important for us, and I'm not going to review it because you definitely learned it in school. At least I think you did. I think I went to school, but it's been a long time since I remember, but...

So, is polynomial long division \inlinemetanote{writes: Poly long div}.

It says the following. So, suppose... and the twist here is it's now for any field, but the fact of the matter is whatever you learned in school also applies here. So, suppose we have two polynomials $f$ and $g$ in this polynomial ring \inlinemetanote{writes: $f(x), g(x) \in K[x]$}. So, coefficients in a field, that's important.

And we're going to do division, so when division, you always have to assume somebody's not zero. This one's not zero \inlinemetanote{writes: $g \neq 0$}. You can't divide by zero.

Then polynomial long division says that there exist unique polynomials \inlinemetanote{writes: $\exists!$} $q(x)$ and $r(x)$ \inlinemetanote{writes: $q(x), r(x)$} that satisfy these two properties. The first property is that $f(x)$ --- $q$ is for quotient, it's kind of the division --- times $g$, and $r$ is the remainder \inlinemetanote{writes: $f(x) = q(x) g(x) + r(x)$}. And then there's the constraint on $r$. So, the constraint is either $r$ is zero \inlinemetanote{writes: either $r(x) = 0$}... and it's the either here because then I'm going to talk about the degree of $r$. I cannot talk about the degree of $r$ formally if it's zero. So, either $r$ is zero or $r$ is not zero \inlinemetanote{writes: or $r(x) \neq 0$} and the degree of $r$ is strictly less than the degree of $g$ \inlinemetanote{writes: and $\deg(r) < \deg(g)$}. Okay, that's the one we're dividing. And you see, I assumed $g$ is not zero, so, I'm allowed to talk about the degree of $g$.
\end{speech}

\begin{content}[Polynomial Long Division (Division with Remainder)]
\begin{theorem}[{Polynomial Long Division in $K[x]$}]\label[theorem]{thm:poly-long-division}
Let $K$ be a field, and let $f(x), g(x) \in K[x]$ with $g(x) \neq 0$.
Then there exist uniquely determined polynomials $q(x), r(x) \in K[x]$ such that:
\begin{equation}\label{eq:poly-long-division}
f(x) = q(x) g(x) + r(x),
\end{equation}
where the remainder $r(x)$ satisfies the dichotomy:
\begin{equation}\label{eq:poly-remainder-condition}
r(x) = 0 \quad \text{or} \quad \bigl( r(x) \neq 0 \ \text{ and } \ \deg(r) < \deg(g) \bigr).
\end{equation}
The polynomial $q(x)$ is called the \newterm{quotient}, and $r(x)$ is called the \newterm{remainder}.
\end{theorem}

\begin{steps}[Handling the Remainder Condition with Undefined Degree]
Because $\deg(0)$ is undefined, one cannot simply state the degree condition as $\deg(r) < \deg(g)$ without excluding the case $r(x) = 0$. The explicit case separation \qt{either $r(x) = 0$ or $\deg(r) < \deg(g)$} preserves strict mathematical consistency with the definition of degree.
\end{steps}
\end{content}

\begin{hidden}
% timestamp: 00:11:00
% topic: Uniqueness of polynomial long division
% board_state: thm:poly-long-division, eq:poly-long-division, eq:poly-remainder-condition
% next_goal: Prove uniqueness by taking the difference of two representations
% open_loops: none
\end{hidden}

\begin{speech}[00:29:16 - 00:30:46]
So, that's polynomial long division, and as I said, you know this. I'm not going to talk about how to prove it. But maybe we talk about some things about the uniqueness. So, polynomial long division is maybe you have some like $3x + 1$, that's your $g$, and this is... I can't believe I'm writing this thing, $x^4$ like this \inlinemetanote{writes: $3x+1 \overline{\smash{\big)}\, x^4 + \dots}$}, and then you divide it, okay? And when you divide it, you get somebody, that's this, and what's left over is that. And if you think about it, they'll satisfy these properties. That's as much as I can possibly say about this.

Okay, so, what about... Did you write it like this?

Okay, good. As I said, I think I went to school, but it's been a long time.

So, maybe something that you didn't emphasize in school is the uniqueness, so, maybe we talk a little bit about that.

So, this says that if I have... \inlinemetanote{writes \qt{$f(x) = q(x) g(x) + r(x)$} and below \qt{$\tilde{q}(x) g(x) + \tilde{r}(x)$}}
\end{speech}

\begin{speech}[00:30:46 - 00:32:21]
I think if you remember every single thing you learned in school, you could almost be a professor.

So, suppose if I have two of these things here. So, if I have two ways to write it, I have to show that they're the same as long as they satisfy these conditions. So, of course, what do I do? I subtract them \inlinemetanote{writes: $0 = (q - \hat{q}) g + (r - \hat{r})$}. I get zero is equal to $q$ minus $q$ hat --- I don't always write the $x$'s, that's just, you know, for emphasis it's a polynomial ---: $(q - \hat{q})g + (r - \hat{r})$. Okay.

So, now then the cases are: so, $g$ is not zero, that was the very beginning. So, suppose this is not zero \inlinemetanote{points to $q - \hat{q}$}. If this is not zero, then this product will have degree at least $g$ \inlinemetanote{writes \qt{$\deg \ge \deg(g)$} under $(q-\hat{q})g$}. And this thing has... again there's some cases about somebody's being zero or not, but this thing can never have degree $g$ or higher, because they're both either zero or born lower. So, that means that the first one has to be zero. So, the first part of this logic is that $q - \hat{q}$ has to be zero \inlinemetanote{writes: $q - \hat{q} = 0$}, which of course means $q = \hat{q}$. Now, once this is zero, then this part dies, then you have this equation, and then again there's some cases, but if they're both not zero, it means they're equal; if one of them is zero, the other one is zero. So, that's the uniqueness. Okay?

All right. \inlinemetanote{erases the long division example on the blackboard}
\end{speech}

\begin{content}[Uniqueness of Polynomial Long Division]
\begin{proof}[Proof of Uniqueness for \cref{thm:poly-long-division}]
Suppose there exist two representations for $f(x)$ upon division by $g(x) \neq 0$:
\begin{align}
f(x) &= q(x) g(x) + r(x), \label{eq:rep1} \\
f(x) &= \tilde{q}(x) g(x) + \tilde{r}(x), \label{eq:rep2}
\end{align}
where both remainder polynomials $r(x)$ and $\tilde{r}(x)$ satisfy condition \eqref{eq:poly-remainder-condition}.
Subtracting \eqref{eq:rep2} from \eqref{eq:rep1} gives:
\begin{equation}\label{eq:uniqueness-diff}
0 = \bigl(q(x) - \tilde{q}(x)\bigr) g(x) + \bigl(r(x) - \tilde{r}(x)\bigr) \implies \bigl(q(x) - \tilde{q}(x)\bigr) g(x) = \tilde{r}(x) - r(x).
\end{equation}
We show that $q(x) - \tilde{q}(x) = 0$:
\begin{description}
    \item[Case $q \neq \tilde{q}$:] Assume for contradiction that $q(x) - \tilde{q}(x) \neq 0$. Since $K$ is a field and $g(x) \neq 0$, the degree of a product equals the sum of the degrees:
    \[
    \deg\Bigl(\bigl(q(x) - \tilde{q}(x)\bigr) g(x)\Bigr) = \deg(q - \tilde{q}) + \deg(g) \ge \deg(g).
    \]
    \item[Degree Contradiction:] On the other hand, each of $r(x)$ and $\tilde{r}(x)$ is either zero or has degree strictly less than $\deg(g)$. The difference $\tilde{r}(x) - r(x)$ is therefore either zero or satisfies:
    \[
    \deg(\tilde{r} - r) \le \max\bigl(\deg(\tilde{r}), \, \deg(r)\bigr) < \deg(g).
    \]
\end{description}
This forces $\deg(g) \le \deg\bigl((q - \tilde{q})g\bigr) = \deg(\tilde{r} - r) < \deg(g)$, an impossible contradiction.
Therefore:
\[
q(x) - \tilde{q}(x) = 0 \implies q(x) = \tilde{q}(x).
\]
Substituting this back into \eqref{eq:uniqueness-diff} yields:
\[
0 \cdot g(x) = \tilde{r}(x) - r(x) \implies \tilde{r}(x) - r(x) = 0 \implies r(x) = \tilde{r}(x).
\]
Hence, the quotient $q(x)$ and remainder $r(x)$ are uniquely determined.
\end{proof}
\end{content}

\begin{aiNote}[Beyond the Lecture: Existence Proof for Polynomial Long Division]
% [PERSISTENT, ADDITION]
The existence of $q(x)$ and $r(x)$ is proved by induction on the degree of $f(x)$. If $\deg(f) < \deg(g)$ or $f = 0$, we set $q(x) = 0$ and $r(x) = f(x)$. If $\deg(f) = n \ge m = \deg(g)$, write $f(x) = a_n x^n + \dots$ and $g(x) = b_m x^m + \dots$. Since $K$ is a field, $b_m \in K^\times$ is invertible. The polynomial $f_1(x) := f(x) - a_n b_m^{-1} x^{n-m} g(x)$ has degree strictly less than $n$ (or is zero). By the induction hypothesis, $f_1(x) = q_1(x) g(x) + r(x)$ with $\deg(r) < \deg(g)$ or $r = 0$, which yields $f(x) = \bigl(a_n b_m^{-1} x^{n-m} + q_1(x)\bigr) g(x) + r(x)$.
\end{aiNote}

\subsection{Every Ideal in \texorpdfstring{$K[x]$}{K[x]} is Principal}

\begin{speech}[00:32:21 - 00:34:09]
But this is an extremely important way to think about these polynomials. I mean, the long division is just an important thing. It sounds silly, but it's important.

It seems like if you went to high school in Basel, you should learn about the Bernoulli numbers, though. If I were in charge of the curriculum... \inlinemetanote{erases the middle chalkboard}

Okay, so, let us say now some things maybe you haven't thought about. So, let $I$ be an ideal \inlinemetanote{writes: Let $I \subset K[x]$ be an ideal}, okay? So, any ideal. Then I claim, so, the proposition --- and this is totally fundamental about polynomial rings --- is that there exists a polynomial in $I$ \inlinemetanote{writes: Prop. $\exists f \in I$ such} such that the whole ideal is just generated by that single thing \inlinemetanote{writes: that $I = (f)$}. It's just like the integers. Integers also have this property: if I take any ideal, it's always just the ideal generated by one person.

So, this is called a principal ideal \inlinemetanote{writes: Principal ideal}. A principal ideal is an ideal generated by one element, and this says that every ideal is principal, just like for the integers. There'll be a definition in a moment. So, let's prove this.
\end{speech}

\begin{content}
\begin{proposition}[{Principal Ideals in $K[x]$}]\label[proposition]{prop:kx-ideals-principal}
Let $K$ be a field, and let $I \subseteq K[x]$ be an ideal. Then there exists a polynomial $f \in I$ such that:
\begin{equation}\label{eq:ideal-generated-single}
I = (f) = f \cdot K[x] = \{ q(x) f(x) \mid q(x) \in K[x] \}.
\end{equation}
That is, every ideal in the polynomial ring $K[x]$ is a \newterm{principal ideal}.
\end{proposition}
\end{content}

\begin{proofWrapper}[Proof of \cref{prop:kx-ideals-principal}]

\begin{speech}[00:34:09 - 00:35:48]
Proof. \inlinemetanote{writes: Proof} So, again, we take $I$ \inlinemetanote{writes: $I \subset K[x]$}. We know nothing about $I$ except it's an ideal.

So, the first case, which is annoying, is that $I$ could be the zero ideal \inlinemetanote{writes: $I = (0)$}. And then we're okay, we're happy \inlinemetanote{writes: $\implies$ happy, $0 \in I$}, since zero is in $I$, because $I$ is... So, this case is fine.

So, the second case is there exist elements $g$ in $I$ such that $g$ is not zero \inlinemetanote{writes: $\exists g \in I, g \neq 0$}. So, either $I$ is just the zero ideal, or there exist some non-zero elements. And now those non-zero elements, we can take their degrees.

So, then, now here's the whole choice, here's the whole proof, actually, in this line: pick $f$ in $I$, $f$ non-zero, such that the degree of $f$ is minimal \inlinemetanote{writes: Pick $f \in I, f \neq 0$ s.t. $\deg(f)$ is minimal}.

Because the set of degrees are some set of integers. As long as it's not empty, well, we can pick a minimal one, and that's what we do. Okay, that's actually the whole proof, this line. You look in $I$, you pick an element such that its degree is minimal, meaning it's minimal among all polynomials in $I$. Maybe I write that: minimal among all elements not zero of $I$ \inlinemetanote{writes: (among all elts $\neq 0$ of $I$)}. All right? So, do you agree I can do that? That has to do with the fact that you can pick lowest integers in a set of non-negative integers. But that's the whole proof here.
\end{speech}

\begin{hidden}
% timestamp: 00:16:30
% topic: Proof that every ideal in K[x] is principal
% board_state: prop:kx-ideals-principal, selection of minimal degree f in I
% next_goal: Divide an arbitrary g in I by f and show the remainder is zero
% open_loops: none
\end{hidden}

\begin{content}[Case 1 and the Choice of the Minimal Degree Generator]
\textbf{Case 1 ($I = \{0\}$):}
If $I$ contains only the zero polynomial, then $I = (0) = 0 \cdot K[x]$, which is principal.

\textbf{Case 2 ($I \neq \{0\}$):}
If $I$ contains non-zero elements, consider the non-empty set of non-negative integer degrees:
\[
\mathcal{D} := \{ \deg(g) \mid g \in I \setminus \{0\} \} \subseteq \mathbb{N}_0.
\]
By the Well-Ordering Principle of $\mathbb{N}_0$, every non-empty set of non-negative integers has a least element. Let $d_{\min} := \min \mathcal{D}$, and choose an element $f \in I \setminus \{0\}$ such that:
\begin{equation}\label{eq:minimal-degree-f}
\deg(f) = d_{\min} = \min \{ \deg(g) \mid g \in I, \, g \neq 0 \}.
\end{equation}
\end{content}

\begin{speech}[00:35:48 - 00:37:19]
So, now we have to prove it. So, then let $g$ in $I$ be any element \inlinemetanote{writes: Let $g \in I$ be any element}. I'm going to use now polynomial long division. Except I'm going to divide $g$ by $f$, so, it's $q(x) \cdot f + r$ \inlinemetanote{writes: $g = q(x) \cdot f + r(x)$}, which is the opposite I did there, but nevertheless, okay?

So, I take this $g$, here's my $f$ that's the lowest, that's the minimal one. I take any $g$ and I use this polynomial long division, all right? So, now remember the cases are... So, the first thing I notice about this polynomial long division is that $f$ is in $I$ \inlinemetanote{writes \qt{$\in I$} under $f$}, so, this whole product is in $I$ \inlinemetanote{writes \qt{$\in I$} under $q(x)f$}, and $g$ is in $I$ \inlinemetanote{writes \qt{$\in I$} under $g$}, so, the conclusion is $r(x)$ is in $I$ \inlinemetanote{writes: $\implies r(x) \in I$}, okay?

For this polynomial long division, I spend my time, eventually get this remainder; the remainder is in $I$, it has to be.
\end{speech}

\begin{content}[Division of an Arbitrary Element by the Minimal Polynomial]
Let $g \in I$ be an arbitrary element. We apply Polynomial Long Division (\cref{thm:poly-long-division}) to divide $g(x)$ by $f(x) \neq 0$:
\begin{equation}\label{eq:div-g-by-f}
g(x) = q(x) f(x) + r(x), \quad \text{where } r(x) = 0 \text{ or } \deg(r) < \deg(f).
\end{equation}
Since $f \in I$ and $I$ is an ideal, the multiple $q(x) f(x) \in I$. Furthermore, $g \in I$.
Rearranging for the remainder yields:
\begin{equation}\label{eq:remainder-in-ideal}
r(x) = \underbrace{g(x)}_{\in I} - \underbrace{q(x) f(x)}_{\in I} \in I.
\end{equation}
Thus, the remainder $r(x)$ is itself an element of the ideal $I$.
\end{content}

\begin{speech}[00:37:19 - 00:38:47]
Okay, so, what are the possibilities for the remainder? The remainder could be zero \inlinemetanote{writes: $r(x) = 0$}. That's good for us, because then again we're happy again, because then that means that $g$ is in \inlinemetanote{writes: $\implies g \in (f)$}... because if this is zero, that means $g$ is in the ideal generated by $f$. Okay?

What's the other possibility? This $r(x)$ is not zero \inlinemetanote{writes: $r(x) \neq 0$} and the degree of $r$ is strictly less than the degree of $f$ \inlinemetanote{writes: $\implies \deg(r) < \deg(f)$}. But that's not possible, because that $r$ would then be an element of $I$ with degree strictly less than $f$, which doesn't exist, because $f$ is minimal with respect to degree.

So, this is not possible \inlinemetanote{writes: not possible} because of the choice of $f$ \inlinemetanote{writes: bec of choice of $f$}.

So, this thing can't happen. We can cross it out with purple \inlinemetanote{crosses out with purple chalk}. So, it has to be this, and that's the proof.
\end{speech}

\begin{content}[Full Restatement: Proof that Every Ideal in $K[x]$ is Principal]
\begin{proof}[Complete Proof of \cref{prop:kx-ideals-principal}]
Let $I \subseteq K[x]$ be an ideal.
\begin{itemize}
    \item \textbf{Zero Ideal:} If $I = \{0\}$, then $I = (0)$ is principal.
    \item \textbf{Minimal Degree Choice:} If $I \neq \{0\}$, choose $f \in I \setminus \{0\}$ of minimal degree among all non-zero polynomials in $I$.
    \item \textbf{First Inclusion ($(f) \subseteq I$):} Clearly, the principal ideal $(f) = f \cdot K[x] \subseteq I$, because $f \in I$ and $I$ is an ideal.
    \item \textbf{Reverse Inclusion ($I \subseteq (f)$):} Conversely, take any $g \in I$. By \cref{thm:poly-long-division}, there exist $q(x), r(x) \in K[x]$ such that:
    \[
    g = q f + r, \quad \text{with } r = 0 \text{ or } \deg(r) < \deg(f).
    \]
    Then $r = g - q f \in I$. If $r \neq 0$, then $r \in I \setminus \{0\}$ would satisfy $\deg(r) < \deg(f)$, directly contradicting the minimality of $\deg(f)$.
    \item \textbf{Conclusion:} Therefore, $r = 0$, which implies:
    \[
    g = q f \in (f).
    \]
    Since $g \in I$ was arbitrary, $I \subseteq (f)$.
\end{itemize}
Combining inclusions gives $I = (f)$, proving that $I$ is principal.
\end{proof}

\begin{steps}[Parallel with the Euclidean Algorithm in $\mathbb{Z}$]
This proof mirrors the classical proof that every ideal in $\mathbb{Z}$ is principal ($I = (n)$):
\begin{itemize}
    \item \textbf{Ring of Integers $\mathbb{Z}$:} One selects the non-zero element $n \in I$ with the smallest absolute value (or smallest positive integer) and uses integer division with remainder $a = q n + r$ ($0 \le r < |n|$).
    \item \textbf{Polynomial Ring $K[x]$:} One selects the non-zero polynomial $f \in I$ with the smallest degree and uses polynomial division with remainder $g = q f + r$ ($r = 0$ or $\deg(r) < \deg(f)$).
\end{itemize}
\end{steps}
\end{content}

\end{proofWrapper}

\subsection{Principal Ideal Domains (PIDs)}

\begin{speech}[00:38:47 - 00:40:01]
So, the outcome of this is that one can think --- of course, this breaks down at a certain point, but it's useful to think about these polynomial rings a little bit like the way you think about the integers. Integers have... I mean, this kind of polynomial in one variable with a field, that they have this property that all the ideals are just generated by one element, just like in the integers. And there's a name for such rings, and they'll play a role later in the class, but maybe I say it.

So, in words, the definition is: $R$ is... so the ring $R$ is a principal ideal domain, which is called PID, but not \qt{pid}. Yeah.
\end{speech}

\begin{student}[Student Question]
Why is $q$ in the ideal?
\end{student}

\begin{speech}[00:40:04 - 00:40:09]
$q$? I never said $q$'s in the ideal. So, I can't answer that question.
\end{speech}

\begin{student}[Student Clarification]
Ah, das ist ein Elementzeichen... \inlinemetanote{realizes the chalk inscription was $f \in I$, not $q \in I$}
\end{student}

\begin{speech}[00:40:09 - 00:40:29]
$f$ is in the ideal, and this is for the product.
\end{speech}

\begin{student}[Student Acknowledgment]
Oh, okay. Thank you.
\end{student}

\begin{speech}[00:40:29 - 00:41:37]
Okay, so, $R$ is a principal ideal domain in words if every ideal $I \subset R$ is principal \inlinemetanote{writes: $R$ is a PID if every ideal $I \subset R$ is principal}.

So, that's in words. It means exactly: for every ideal in $R$, there exists some element of the ideal such that the ideal is generated by that element \inlinemetanote{writes: $\left[\forall \text{ ideal } I \subset R, \ \exists f \in I \text{ s.t. } I = (f)\right]$}, okay? This is just repeating it.

So, that's a special kind of domain. I should say that when I say $R$ is a principal ideal domain, it should be a domain too. I didn't write it, I guess, because it's already there. $R$ is a principal ideal domain if every ideal of $R$ is principal, and $R$ is a domain \inlinemetanote{adds: and $R$ is a domain}.

Okay? But it's already in the... right. So, to be a principal ideal domain, you have to be a domain, of course.

So, what we've seen in the class so far: $\mathbb{Z}$, I mean, many of them, but let's say $\mathbb{Z}$ is a PID \inlinemetanote{writes: $\mathbb{Z} = \text{PID}$}.

What about $\mathbb{Z} \bmod 6\mathbb{Z}$? \inlinemetanote{writes: $\faktor{\mathbb{Z}}{6\mathbb{Z}}$} Well, no, because it's not a domain \inlinemetanote{writes: $\text{not a domain}$}.

A field is a principal ideal domain. It's kind of a dumb principal ideal domain, because there's not many ideals, but it is \inlinemetanote{writes: A field is a PID, because $(0), (1)$}. Because there are only two ideals in a field: one ideal is this one, and one ideal is that one. So, it's also a principal ideal domain --- not, as I said, not the most interesting way to look at a field. That's an interesting principal ideal domain, and then today we see that if $K$ is a field, then this is a principal ideal domain \inlinemetanote{writes: If $K$ is a field $\implies K[x] = \text{PID}$}.
\end{speech}

\begin{hidden}
% timestamp: 00:22:00
% topic: Definition and examples of Principal Ideal Domains (PIDs)
% board_state: def:pid, ex:pid-examples, non-examples Z[x] and C[x, y]
% next_goal: Discuss non-PIDs with multiple generators and wrap up the lecture
% open_loops: none
\end{hidden}

\begin{content}
\begin{definition}[Principal Ideal Domain (PID)]\label[definition]{def:pid}
An integral domain $R$ is called a \newterm{principal ideal domain} (abbreviated \newterm{PID}) if every ideal $I \subseteq R$ is principal, meaning:
\begin{equation}\label{eq:pid-condition}
\forall \text{ ideal } I \subseteq R, \ \exists f \in I \quad \text{such that} \quad I = (f).
\end{equation}
\end{definition}

\begin{examples}[Examples and Non-Examples of PIDs]\label[examples]{ex:pid-overview}
\begin{enumerate}[label=\textbf{(\arabic*)}]
    \setcounter{enumi}{0} \item \textbf{The Integers $\mathbb{Z}$ (\textcolor{ForestGreen}{PID}):} Every ideal in $\mathbb{Z}$ has the form $n\mathbb{Z} = (n)$. Since $\mathbb{Z}$ is an integral domain, $\mathbb{Z}$ is a PID.
    \setcounter{enumi}{1} \item \textbf{Quotient Ring \texorpdfstring{$\faktor{\mathbb{Z}}{6\mathbb{Z}}$}{Z/6Z} (\textcolor{BrickRed}{NOT a PID}):} Although every ideal in $\faktor{\mathbb{Z}}{6\mathbb{Z}}$ is principal, $\faktor{\mathbb{Z}}{6\mathbb{Z}}$ is \emph{not} an integral domain because it contains zero divisors ($2 \cdot 3 \equiv 0 \pmod 6$).
    \setcounter{enumi}{2} \item \textbf{Any Field $K$ (\textcolor{ForestGreen}{PID}):} If $K$ is a field, its only ideals are $(0)$ and $(1) = K$. Both are principal, and $K$ is an integral domain, so every field is a PID.
    \setcounter{enumi}{3} \item \textbf{Polynomial Ring $K[x]$ over a Field (\textcolor{ForestGreen}{PID}):} By \cref{prop:kx-ideals-principal}, every ideal in $K[x]$ is principal. Since $K[x]$ has no zero divisors when $K$ is a field, $K[x]$ is a PID.
\end{enumerate}
\end{examples}
\end{content}

\subsection{Rings that are Not Principal Ideal Domains}

\begin{speech}[00:41:37 - 00:41:50]
Something I should have said kind of clearly here is that that division algorithm that you learn in school uses the fact that you have a field, because, you know, you might have to divide by the leading coefficients. So, that was used in the way that the division algorithm works.

So, some things that are not... rings that are not principal ideal domains \inlinemetanote{writes: Rings that are not PIDs}.

Well, one which is kind of close is that if I take polynomials over $\mathbb{Z}$ in one variable \inlinemetanote{writes: $\mathbb{Z}[x]$}, because you can't use the division algorithm the same way, because you can't divide. So, can you give me an ideal here that's probably not principal?

I don't expect you to give me a proof of it, but what do you think would be probably not principal here?
\end{speech}

\begin{student}[Student Suggestion]
$x^2 + 1$?
\end{student}

\begin{speech}[00:41:52 - 00:42:03]
Yeah, I think that might be principal. But, you know, if you're just using the $x$ part with leading terms $1$... yeah, okay, you can think about that ideal, I think that one is principal. The standard example is you use something from $\mathbb{Z}$, like the $2$, and something from $x$ on this side \inlinemetanote{writes: $\supset (2, x)$}. This is like the canonical one.

Yeah, yours is principal, I'd say. I think.

All right, so, this is the standard one that's not principal. And the reason, as I said, you take something from the $\mathbb{Z}$ side and something from the two side. And so, this should be an exercise for you: prove it's not principal \inlinemetanote{writes: Prove not principal}.

And another way you could do it, if you have a field with polynomials in two variables over a field, then you can take the ideal $(x, y)$ \inlinemetanote{writes: $\mathbb{C}[x, y] \supset (x, y)$} \inlinemetanote{writes: Prove not princ.}. And that's also clearly not going to be principal, because there's no way you can find one polynomial that generates both $x$ and $y$. Okay. That's it. \inlinemetanote{students applaud}
\end{speech}

\begin{content}[Non-Principal Ideals in Multi-Variable and Integer Polynomial Rings]
\begin{examples}[Rings that are Not Principal Ideal Domains]\label[examples]{ex:non-pids}
\begin{enumerate}[label=\textbf{(\arabic*)}]
    \setcounter{enumi}{0} \item \textbf{Polynomials over the Integers $\mathbb{Z}[x]$ (\textcolor{BrickRed}{NOT a PID}):}
    The ideal generated by $2$ and $x$:
    \begin{equation}\label{eq:ideal-2-x}
    I := (2, x) = \{ 2 p(x) + x q(x) \mid p(x), q(x) \in \mathbb{Z}[x] \}
    \end{equation}
    consists of all polynomials in $\mathbb{Z}[x]$ whose constant term is even. It cannot be generated by a single polynomial, so $\mathbb{Z}[x]$ is not a PID.
    
    \setcounter{enumi}{1} \item \textbf{Polynomials in Two Variables $\mathbb{C}[x, y]$ (\textcolor{BrickRed}{NOT a PID}):}
    In the polynomial ring in two variables over $\mathbb{C}$ (or any field $K$), the maximal ideal:
    \begin{equation}\label{eq:ideal-x-y}
    J := (x, y) = \{ x p(x, y) + y q(x, y) \mid p, q \in \mathbb{C}[x, y] \}
    \end{equation}
    consisting of all polynomials without constant term, requires at least two generators. No single polynomial in $\mathbb{C}[x, y]$ can divide both $x$ and $y$ simultaneously while generating all elements in $J$; a proof is given at the start of Week~6 (\cref{ans:kx-ky-pid}).
\end{enumerate}
\end{examples}

\begin{exercise}[{Proving that $(2, x)$ is Not Principal in $\mathbb{Z}[x]$}]\label[exercise]{exer:prove-2-x-not-principal}
Prove that the ideal $(2, x) \subset \mathbb{Z}[x]$ is not principal.
\end{exercise}

\begin{steps}[Proof Hint for \cref{exer:prove-2-x-not-principal}]
Suppose $(2, x) = (f(x))$ for some $f(x) \in \mathbb{Z}[x]$. Since $2 \in (f(x))$, $f(x)$ must divide $2$ in $\mathbb{Z}[x]$, which forces $f(x) = \pm 1$ or $f(x) = \pm 2$. If $f(x) = \pm 1$, then $(f(x)) = \mathbb{Z}[x]$, but $(2, x)$ contains only polynomials with even constant terms, so $1 \notin (2, x)$. If $f(x) = \pm 2$, then $x \in (2)$, which would mean $x = 2 q(x)$ for some $q(x) \in \mathbb{Z}[x]$, which is impossible since the leading coefficient of $x$ is $1$, not divisible by $2$. Thus, no single generator exists.
\end{steps}
\end{content}

\begin{overview}[Summary of Key Concepts]
\begin{itemize}
    \item \textbf{Canonical Product Homomorphism:} For any ideals $I, J \subseteq R$, the natural map $\phi \colon \faktor{R}{IJ} \to \faktor{R}{I} \times \faktor{R}{J}$ given by $\bar{z} \mapsto (\bar{z}, \bar{z})$ is a well-defined ring homomorphism. Without comaximality, $\phi$ generally fails to be injective or surjective (e.g., $I = (4)$ and $J = (6)$ in $\mathbb{Z}$).
    \item \textbf{Chinese Remainder Theorem:} When $I + J = (1) = R$, there exist $i \in I$ and $j \in J$ such that $i + j = 1$. This partition of unity ensures that $\phi$ is an isomorphism, yielding explicit solutions $z = jx + iy$ to simultaneous systems of congruences.
    \item \textbf{Arithmetic in $\mathbb{Z}$ and Euler's Totient Function:} Iterated application yields $\faktor{\mathbb{Z}}{n\mathbb{Z}} \cong \prod_{i=1}^r \faktor{\mathbb{Z}}{p_i^{k_i}\mathbb{Z}}$. Taking unit groups gives $\varphi(n) = n \prod_{p \mid n} (1 - 1/p)$. The function mixes primes via $p - 1$, preventing direct inversion (a cornerstone of cryptography and connected to irregular primes such as $691 \mid \text{numerator}(B_{12})$).
    \item \textbf{Polynomial Rings $K[x]$ over Fields:} For a field $K$, polynomials enjoy Euclidean long division with unique quotient and remainder ($r = 0$ or $\deg(r) < \deg(g)$). The degree of the zero polynomial $\deg(0)$ is left strictly undefined.
    \item \textbf{Principal Ideal Domains (PIDs):} Using the well-ordering of degrees, every ideal $I \subseteq K[x]$ is generated by any non-zero element of minimal degree, proving that $K[x]$ is a PID. Multi-variable polynomial rings like $\mathbb{C}[x, y]$ and polynomial rings over non-fields like $\mathbb{Z}[x]$ fail to be PIDs due to non-principal ideals such as $(x, y)$ and $(2, x)$.
\end{itemize}
\end{overview}